\documentclass[a4paper, 10pt]{article}
\usepackage[top=2.0cm,bottom=2.0cm,left=1.5cm,right=2.5cm]{geometry}

\usepackage{float}
\usepackage{hyperref}
\hypersetup{
  colorlinks   = true,   % Colours links instead of ugly boxes
  urlcolor     = red,    % Colour for external hyperlinks
  linkcolor    = red,    % Colour of internal links
  citecolor    = red     % Colour of citations
}

\usepackage[ngerman]{babel}
\usepackage[ansinew]{inputenc}

% Notes am rechten Seitenrand -deshalb die unterschiedliche Breite.
\usepackage{marginnote}

\usepackage{amssymb}
\usepackage{amsmath}
\usepackage{amstext}
\usepackage{amsfonts}
\usepackage{eurosym}
\usepackage{trfsigns}

\usepackage{longtable}
\usepackage{tabularx}
\usepackage{multirow}
\usepackage[table]{xcolor}
\usepackage{graphicx}
\usepackage{booktabs}% http://ctan.org/pkg/booktabs
\newcommand{\tabitem}{~~\llap{\textbullet}~~}

\usepackage{tcolorbox}
\usepackage{pdflscape}

\tcbuselibrary{breakable,skins}
\usepackage{ctable}
\definecolor{violett}{RGB}{154,50,205 }
\definecolor{darkgreen}{RGB}{0,136,0 }
\definecolor{htwgreen}{RGB}{134,189,36 }

% alternate rowcolors for all long-tables
\let\oldlongtable\longtable
\let\endoldlongtable\endlongtable
\renewenvironment{longtable}{\rowcolors{2}{white}{lightgray}\oldlongtable} {
\endoldlongtable}

% Definition
% #1: Title
% #2: Bemerkung
\newcounter{def}
\newtcolorbox{defbox}[3][]{
  enhanced,
  overlay unbroken and first={ 
\node[fill=white,rounded corners, draw=white, text=gray, line width=3pt, inner sep=4pt, anchor=east, xshift=-10pt] at (frame.north east){#2};
\node[anchor=north east, text=gray, fill=white, text width=3.5em, align=left, minimum height=4ex] at ([xshift=4.5\marginparsep]frame.north east) {\small Def. \stepcounter{def}\thedef \newline \textcolor{black}{#3} };
}, 
  colframe=gray, 
  %colback=mytestcolor, 
  colback=white, 
  coltitle=red!70!gray, 
  boxrule=1pt, 
  arc=1mm, 
  breakable,
  top=5pt,
  before=\vskip18pt}

% Satz
% #1: Title
% #2: Bemerkung
\newcounter{satz}
   \newtcolorbox{theorembox}[3][]{
  enhanced,
  overlay unbroken and first={ 
\node[fill=white,rounded corners, draw=white, text=gray, line width=3pt, inner sep=4pt, anchor=east, xshift=-10pt] at (frame.north east){#2};
\node[anchor=north east, text=gray, fill=white, text width=3.5em, align=right, minimum height=4ex] at ([xshift=4.5\marginparsep]frame.north east) {\small Regel \stepcounter{satz}\thesatz \newline \textcolor{black}{#3} };
}, 
  colframe=gray, 
  %colback=mytestcolor, 
  colback=white, 
  coltitle=red!70!gray, 
  boxrule=1pt, 
  arc=1mm, 
  breakable,
  top=5pt,
  before=\vskip18pt
}

% Beispiel
% #1: Title
% #2: Bemerkung
\newcounter{exa}
\newtcolorbox{exabox}[3][]{
  enhanced,
  overlay unbroken and first={ 
	  \node[fill=white, rounded corners, draw=white, text=blue, line width=3pt, inner sep=4pt, anchor=east, xshift=-10pt] at (frame.north east) {#2}; 
	  \node[anchor=north east, text=blue, fill=white, text width=3.5em, align=right, minimum height=4ex] at ([xshift=4.5\marginparsep]frame.north east)  {\small  Bsp \stepcounter{exa}\theexa \newline \textcolor{black}{ #3} };
}, 
  colframe=blue, 
  colback=white, 
  coltitle=red!70!blue, 
  boxrule=1pt, 
  arc=1mm, 
  breakable,
  top=5pt,
  before=\vskip18pt
}

% �bung
% #1: Title
% #2: Bemerkung
\newcounter{exercise}
\newtcolorbox{exercisebox}[1][]{
  enhanced,
  overlay unbroken and first={ 
	  \node[fill=white, rounded corners, draw=white, text=black, line width=3pt, inner sep=4pt, anchor=east, xshift=-10pt] at (frame.north east)  {#1}; 
	  \node[anchor=north east, text=black, fill=white, text width=3.5em, align=right, minimum height=4ex] at ([xshift=4.5\marginparsep]frame.north east)  
	  {\small  Aufg. \stepcounter{exercise}\theexercise };
}, 
  colframe=black, 
  colback=white, 
  coltitle=red!70!blue, 
  boxrule=1pt, 
  arc=1mm, 
  breakable,
  top=5pt,
  before=\vskip18pt
}

% Aufgaben gemeinsam mit der Theorie verlinken; siehe LaTeX/README.md.
% Einmal in der Praeambel des Hauptdokuments nach hyperref laden.
% true: Theorie, Aufgaben und Loesungen gemeinsam; false: getrennte Ausgaben.
\newif\ifintegriereaufgaben
\integriereaufgabenfalse

% Fuehrt den gesamten Inhalt nur bei vorhandenem Label aus.
% Wie normale Querverweise benoetigt die Erkennung mehrere LaTeX-Laeufe.
\newcommand{\wennreferenzvorhanden}[2]{%
  \ifcsname r@#1\endcsname
    #2%
  \fi
}

% Verlinkte Nummer bei vorhandenem Ziel, andernfalls keinerlei Ausgabe.
\newcommand{\abschnittsverweis}[1]{%
  \wennreferenzvorhanden{#1}{\ref{#1}}%
}

% Gleichbedeutender Name fuer bestehende Theoriehinweise.
\newcommand{\skriptverweis}[1]{\abschnittsverweis{#1}}

% Ohne Ziel wird auch kein leerer Randnotiz-Kasten angelegt.
\newcommand{\randverweis}[1]{%
  \wennreferenzvorhanden{#1}{\marginpar{\small\abschnittsverweis{#1}}}%
}

\integriereaufgabenfalse

\begin{document}
\pagenumbering{roman}

\noindent
{
\tableofcontents
\newpage

\pagenumbering{arabic}
\setcounter{section}{0}
% Matrizen
\setcounter{satz}{0} \setcounter{def}{0} \setcounter{exercise}{0}  

\ifintegriereaufgaben
\section*{Matrizen Aufgaben}\label{MA1-matrizen-aufgaben}

	Regeln und Beispiele, die zur Bearbeitung der Aufgaben hilfreich sind, finden Sie im Skript\wennreferenzvorhanden{MA1-matrizen}{ (Abschnitt \abschnittsverweis{MA1-matrizen})} und, falls Sie nicht weiterkommen, schauen Sie \hyperref[MA1-matrizen-loesungen]{hier}.
\else 
\section{Matrizen Aufgaben}\label{MA1-matrizen-aufgaben}

	Regeln und Beispiele, die zur Bearbeitung der Aufgaben hilfreich sind, finden Sie im Skript und, falls Sie nicht weiterkommen, schauen Sie \hyperref[MA1-matrizen-loesungen]{hier}.
\fi

\begin{exercisebox}[Rechenregeln]
	Berechnen Sie $ (  A^\intercal C - 2\cdot B)(-3\cdot C^\intercal )$ f�r die folgenden Matrizen
	\[ A = \left( \begin{array}{ccc} 1 & -2 & 3 \\ 0 & 3 & 1 \end{array}\right), \quad B = \left( \begin{array}{cc} 2 & -1 \\ 0 & 3 \\ -2 & -2 \end{array}\right), \quad C=\left( \begin{array}{cc} 3 & 2 \\ -1 & 2 \end{array}\right) \,. \]
\end{exercisebox}

%\begin{exercisebox}[Matrixmultiplikation und Skalarprodukt]
%Es sei $\boldsymbol a \in \mathbb R^n, \boldsymbol b\in \mathbb R^m$ und $A\in \mathbb R^{n\times m}$.  Zeigen Sie $\langle \boldsymbol a , A \boldsymbol b\rangle = \langle A^\intercal \boldsymbol a, \boldsymbol b \rangle\,.$
%\end{exercisebox}

\begin{exercisebox}[Rang und Bild]
	Es sei $A = \boldsymbol a \boldsymbol b^\intercal, \; \boldsymbol a \in \mathbb R^n, \boldsymbol b \in \mathbb R^m$. Wie ist der Rang der Matrix $A$? Wie sieht das Bild der Matrix aus?
\end{exercisebox}

\begin{exercisebox}[Basis]
	\begin{itemize}
		\item[(a)] F�r welche Werte von $a\in\mathbb R$ bilden die vier Vektoren eine Basis des $\mathbb R^4$?
  \begin{eqnarray*}
	  \boldsymbol x_1 = (3,1,4,0)^\intercal, \;  
	  \boldsymbol x_2 = (1,1,0,6)^\intercal, \; 
	  \boldsymbol x_3 = (-4,0,5,a)^\intercal, \;  
	  \boldsymbol x_4 = (0,0,1,2)^\intercal\,.
  \end{eqnarray*}

\item[(b)]	  �berpfr�fen Sie, ob die Polynome $1, 1+x, 1-x, x^2$ eine Basis des $\mathcal P_2$ bilden.
\end{itemize}
\end{exercisebox}

\begin{exercisebox}[Lineare Abbildungen]
	Gegeben sind die linearen Abbildungen 
	\[ \mathcal A: \mathbb R^2 \to \mathbb R^3, \left( \begin{array}{c} x_1 \\ x_2 \end{array}\right) \mapsto \left( \begin{array}{c} 0 \\ 3x_1-x_2 \\ 2x_2 \end{array}\right), \quad \mathcal B:\mathbb R^2 \to \mathbb R^3, \left( \begin{array}{c} x_1 \\ x_2 \end{array}\right) \mapsto \left( \begin{array}{c} -x_1+x_2 \\ -x_2 \\ 3x_2 \end{array}\right)\,.\]
	\begin{itemize}
		\item[(a)] Bestimmen Sie die darstellenden Matrizen $A$ und $B$ der linearen Abbildungen $\mathcal A$ und $\mathcal B$.
		\item[(b)] Gegeben ist eine Matrix \[ C=\left( \begin{array}{ccc} -1 & 2 & 0 \\ 0 & 3 & -1 \end{array}\right) \in \mathbb R^{2\times 3}\] Bestimmen Sie die von $C$ induzierte lineare Abbildung $\mathcal C$.
		\item[(c)] Wie lautet die darstellende Matrix der linearen Abbildung $\mathcal C\circ \mathcal A$?
	\end{itemize}
\end{exercisebox}

\begin{exercisebox}[Lineare Abbildungen]
	Gegeben sei der reelle Vektorraum \( \mathcal P_3:= \left\{ p:\mathbb R\to \mathbb R\;|\;p \;\textnormal{Polynome vom Grad } 3\right\}\) und die Abbildung \( f: \mathcal P_3 \to \mathbb R^2, \quad p\mapsto \left( \begin{array}{c}p(2) \\ p(3)\end{array}\right) \,.\)
	\begin{itemize}
		\item[(a)] Zeigen Sie, dass $f$ eine lineare Abbildung ist und berechnen Sie $f$ f�r die Monome $p_i: x\mapsto x^i,\; i=0,1,2,3\,.$
		\item[(b)] Berechnen Sie $f(p)$ f�r ein allgemeines Polynom $p \in \mathcal P_4$. % mit  \[ p=a_3\cdot p_3 + a_2 \cdot p_2 + a_1\cdot p_1 + a_0 \cdot p_0, \quad a_0,\dots,a_3\in \mathbb R.\] 
		\item[(d)] Zeigen Sie, dass die Funktion \[ g:\mathbb R^4 \to \mathbb R^2, \quad \left( \begin{array}{c}a_3 \\ a_2 \\ a_1 \\ a_0\end{array}\right) \mapsto f(a_3\cdot p_3 + a_2 \cdot p_2 + a_1\cdot p_1 + a_0 \cdot p_0)\] linear ist und berechnen Sie die darstellende Matrix von $g$.
	\end{itemize}
\end{exercisebox}

%\begin{exercisebox}[Mathematik lesen, Eigenwerte]
%\end{exercisebox}
%\begin{exercisebox}[Mathematik lesen, Eigenvektoren]
%\end{exercisebox}
%\begin{exercisebox}[Mathematik lesen, Positive Definitheit]
%\end{exercisebox}
%\begin{exercisebox}[Mathematik lesen, Symmetrie]
%\end{exercisebox}
%\begin{theorembox}[Grafikkarten Programmierung - cuda]
%\end{theorembox}
%\begin{theorembox}[D�nnbesetzte Matrizen]
%\end{theorembox}
%\begin{theorembox}[Hierarchische Matrizen, Niedrigrangapproximation]
%\end{theorembox}
%\begin{theorembox}[Spezielle Matrizen  (Elemente erg�nzen vielleicht?)- Mathematik lesen oder Beispiele angeben?]
%\begin{itemize}
%	\item Orthogonale Matrix
%	\item Komplex konjugierte
%	\item antisymmetrische Matrix
%	\item Singul�re Matrix
%	\item Transponierte
%	\item  $A \in \mathbb C^{2\times 3}$. Wie lautet $A^\intercal$? Konjugierte 
%	\item Basis des Vektorraums $\mathbb C^{2\times 3}$
%\end{itemize}
%\end{theorembox}
 
\newpage
\setcounter{satz}{0} \setcounter{def}{0} \setcounter{exercise}{0}  
\section*{Matrizen L�sungen}\label{MA1-matrizen-loesungen}

\ifintegriereaufgaben
	Regeln und Beispiele, die zum Verst�ndnis der L�sungswege hilfreich sind, finden Sie im Skript\wennreferenzvorhanden{MA1-matrizen}{ (Abschnitt \abschnittsverweis{MA1-matrizen})}.
\else
	Regeln und Beispiele, die zum Verst�ndnis der L�sungswege hilfreich sind, finden Sie im Skript.
\fi 

\begin{exercisebox}[Rechenregeln]
	Gegeben seien die Matrizen 
	\[ A = \left( \begin{array}{ccc} 1 & -2 & 3 \\ 0 & 3 & 1 \end{array}\right), \quad B = \left( \begin{array}{cc} 2 & -1 \\ 0 & 3 \\ -2 & -2 \end{array}\right), \quad C=\left( \begin{array}{cc} 3 & 2 \\ -1 & 2 \end{array}\right) \,. \]
	Berechnen Sie $ ( A^\intercal C - 2\cdot B)(-3\cdot C^\intercal )\,.$

\vspace{0.3cm}
\noindent\textbf{L�sung:}\newline
\[ \left( \begin{array}{cc} -15 & -27 \\ 105 & -3 \\ -180 & -36 \end{array}\right)\,. \]
\end{exercisebox}

%\begin{exercisebox}[Matrixmultiplikation und Skalarprodukt]
%Es sei $\boldsymbol a \in \mathbb R^n, \boldsymbol b\in \mathbb R^m$ und $A\in \mathbb R^{n\times m}$.  Zeigen Sie $\langle \boldsymbol a , A \boldsymbol b\rangle = \langle A^\intercal \boldsymbol a, \boldsymbol b \rangle\,.$
%
%\vspace{0.3cm}
%\noindent\textbf{L�sung:}\newline
%\begin{eqnarray*}
%	\langle \boldsymbol a, A\boldsymbol b\rangle &=& \sum\limits_{i=1}^n \left( a_i \cdot \left( A \boldsymbol b\right)_i \right) \\
%	&=& \sum\limits_{i=1}^n a_i \left( \sum\limits_{j}^m a_{ij} b_j \right) \\
%	&=& \sum\limits_{i=1}^n\sum\limits_{j}^m  a_i a_{ij} b_j  \\
%	&=&  \sum\limits_{j}^m \left(\sum\limits_{i=1}^n  \underbrace{a_i a_{ij}}_{\displaystyle{= \boldsymbol a^\intercal A = \left( \boldsymbol a A^\intercal \right)^\intercal = \boldsymbol a A^\intercal }} \right) b_j  \\
%	&=&  \sum\limits_{j}^m \left(A^\intercal \boldsymbol a \right)_j b_j
%\end{eqnarray*}
%\end{exercisebox}

\begin{exercisebox}[Rang und Bild]
	Es sei $A = \boldsymbol a \boldsymbol b^\intercal, \quad \boldsymbol a \in \mathbb R^n, \boldsymbol b \in \mathbb R^m$. Wie ist der Rang der Matrix $A$? Wie sieht das Bild der Matrix aus?


\vspace{0.3cm}
\noindent\textbf{L�sung:}\newline
Der Rang der Matrix $A \in \mathbb R^{n\times m}$ entspricht der Dimension des Bildes der linearen Abbildung $A:\mathbb R^m \to \mathbb R^n$. Es gilt 
\begin{eqnarray*} 
	\forall \boldsymbol y \in \mathbb R^m: A\boldsymbol y &=& (\boldsymbol a \boldsymbol b^\intercal) \boldsymbol y = \boldsymbol a \underbrace{(\boldsymbol b^\intercal  \boldsymbol y)}_{\displaystyle{ =: \lambda \in \mathbb R!}} = \lambda \boldsymbol a\,,
\end{eqnarray*}
wobei wir das Assoziativgesetz der Matrixmultiplikation, $(A B)C = A(BC)$, ausgenutzt haben. Die Multiplikation eines Zeilenvektors mit einem Spaltenvektor gleicher Gr��e ist eine reelle Zahl\footnote{ im Grunde ist diese Zahl das Skalarprodukt der beiden Vektoren: \(
\boldsymbol b^\intercal  \boldsymbol y = \langle \boldsymbol b ,\boldsymbol y\rangle =:  \lambda \in \mathbb R\,. \)}
Damit umfasst das Bild der linearen Abbildung $A$ alle skalaren Vielfachen des Vektors $\boldsymbol a$ und damit
\[ \mathrm{rk}(A) = \mathrm{dim}\left( \mathrm{im}A\right) = \mathrm{dim} \{ \boldsymbol x \in \mathbb R^n: \boldsymbol x = \lambda \boldsymbol a\} = 1\,.\]
\end{exercisebox}

\begin{exercisebox}[Basis]
	\begin{itemize}
		\item[(a)] F�r welche Werte von $a\in\mathbb R$ bilden die vier Vektoren eine Basis des $\mathbb R^4$?
  \begin{eqnarray*}
	  \boldsymbol x_1 = (3,1,4,0)^\intercal, \;  
	  \boldsymbol x_2 = (1,1,0,6)^\intercal, \; 
	  \boldsymbol x_3 = (-4,0,5,a)^\intercal, \;  
	  \boldsymbol x_4 = (0,0,1,2)^\intercal\,.
  \end{eqnarray*}

\item[(b)]	  �berpr�fen Sie, ob die Polynome $1, 1+x, 1-x, x^2$ eine Basis des $\mathcal P_2$ bilden.
\end{itemize}

  %\hspace{0.3cm}
  %\newline
  \textbf{L�sung:}
 \begin{itemize}
	 \item[(a)] Die Vektoren sind genau dann linear unabh�ngig, wenn die Matrix
		 \[ A_a := \left( \begin{array}{cccc} 3 & 1 & -4 & 0 \\ 1 & 1 & 0 & 0 \\ 4 & 0 & 5 & 1 \\ 0 & 6 & a & 2 \end{array}\right) \in \mathbb R^{4\times 4}\,, \]
		 
		 deren Spalten $\boldsymbol x_1, \boldsymbol x_2, \boldsymbol x_3$ und $\boldsymbol x_4$ sind, vollen Rang hat. Wir bringen $A_a$ auf Dreiecksform und bestimmen den Rang in Abh�ngigkeit von $a\in \mathbb R$. Es ist geschickt die zwei letzten Spalten zu vertauschen, so dass der Parameter \(a\) rechts unten landet und so in m�glichst wenig Rechenoperationen mitgeschleppt wird:      
		 \begin{eqnarray*}
			 \left( \begin{array}{cccc} 3 & 1 & -4 & 0 \\ 1 & 1 & 0 & 0 \\ 4 & 0 & 5 & 1 \\ 0 & 6 & a & 2 \end{array}\right) &\Leftrightarrow&
			 \left( \begin{array}{cccc} 3 & 1 & 0 & -4 \\ 1 & 1 & 0 & 0 \\ 4 & 0 & 1 & 5 \\ 0 & 6 & 2 & a \end{array}\right) \Leftrightarrow 
		 \left( \begin{array}{cccc} 3 & 1 & 0 & -4 \\ 0 & -2 & 0 & -4 \\ 0 & 4 & -3 & -31 \\ 0 & 6 & 2 & a \end{array}\right) \\
		 \left( \begin{array}{cccc} 3 & 1 & 0 & -4 \\ 0 & -2 & 0 & -4 \\ 0 & 0 & -3 & -39 \\ 0 & 0 & 2 & -12+a \end{array}\right) &\Leftrightarrow& 
		 \left( \begin{array}{cccc} 3 & 1 & 0 & -4 \\ 0 & -2 & 0 & -4 \\ 0 & 0 & -3 & -39 \\ 0 & 0 & 0 & -114-3a \end{array}\right) \,.
	 \end{eqnarray*}
	 F�r $a=38$ hat die Matrix $A_{38}$ also Rang $3<4$ beliebig und damit sind die Vektoren f�r $a=38$ linear abh�ngig. F�r $a\not = 38$ hat die Matrix vollen Rang und das hei�t die Vektoren sind linear unabh�ngig.
 \item[(b)]
	 Zwei L�sungsvarianten: 
	 \begin{itemize}
		 \item Die drei Polynome \(\{p_0(x)= 1, p_1(x)= 1+x, P_2(x)= 1-x\}\) sind linear abh�ngig, weil \[ p_0(x) = \frac12 p_1(x) + \frac12 p_2(x)\,.\] Also kann es sich nicht um eine Basis handeln.  
		 \item Der systematische Weg zu pr�fen, ob die gegebenen Polynome linear abh�ngig sind ist folgender: man zeige, dass die Nullabbildung \(x\mapsto 0\) auf nichttriviale Weise darstellbar ist, dass also die Gleichung 
	 \[ \alpha_1 \cdot 1 + \alpha_2 \cdot (1+x) + \alpha_3 \cdot (1-x) +\alpha_4 \cdot x^2 = 0 \quad (*)\,,\]
	 mindestens eine nichttriviale L�sung besitzt. Nichttriviale L�sung hei�t, dass f�r mindestens einen Koeffizienten \(\alpha_i, i\in \{1,2,3,4\},\) gilt $\alpha_i\not= 0$. Aus $(*)$ folgt durch Ausklammern und Zusammenfassen 
	\begin{eqnarray*}
		1\cdot (\alpha_1 + \alpha_2 + \alpha_3) + x\cdot ( \alpha_2 - \alpha_3) + \alpha_4 x^2 = 0 \,,
	\end{eqnarray*}
	und hieraus folgt f�r die unbekannten Koeffizienten $\alpha_i, \; i=1,2,3,4$:
	\begin{eqnarray*}
		\left\{ \begin{array}{lcl}
		\alpha_1 + \alpha_2 + \alpha_3 &=& 0 \\
		\alpha_2 - \alpha_3 &=& 0\\
\alpha_4 &=& 0 
\end{array}\right.
\end{eqnarray*}
Dieses Gleichungssystem ist unterbestimmt und hat unendlich viele L�sungen:
\begin{eqnarray*}
	\Rightarrow \begin{cases}{r}  \alpha_1 = -2t \\ \alpha_2=\alpha_3=t\in \mathbb R  \\ \alpha_4 = 0\,.
	\end{cases}
\end{eqnarray*}
Die Polynome $1, 1+x, 1-x, x^2$ bilden also kein minimales Erzeugendensystem, sind also keine Basis des $\mathcal P_2$. Zum Beispiel erhalten wir f�r $t=1$ die Koeffizienten $\alpha_1 = -2, \alpha_2=1, \alpha_3=1, \alpha_4 =0$, die eine nichttriviale Linearkombination bilden. 
\end{itemize}
\end{itemize}
\end{exercisebox}

\begin{exercisebox}[Lineare Abbildungen]
	Gegeben sind die linearen Abbildungen 
	\[ \mathcal A: \mathbb R^2 \to \mathbb R^3, \left( \begin{array}{c} x_1 \\ x_2 \end{array}\right) \mapsto \left( \begin{array}{c} 0 \\ 3x_1-x_2 \\ 2x_2 \end{array}\right), \quad \mathcal B:\mathbb R^2 \to \mathbb R^3, \left( \begin{array}{c} x_1 \\ x_2 \end{array}\right) \mapsto \left( \begin{array}{c} -x_1+x_2 \\ -x_2 \\ 3x_2 \end{array}\right)\,.\]
	\begin{itemize}
		\item[(a)] Bestimmen Sie die darstellenden Matrizen $A$ und $B$ der linearen Abbildungen $\mathcal A$ und $\mathcal B$.
		\item[(b)] Gegeben ist eine Matrix \[ C=\left( \begin{array}{ccc} -1 & 2 & 0 \\ 0 & 3 & -1 \end{array}\right) \in \mathbb R^{2\times 3}\] Bestimmen SIe die von $C$ induzierte lineare Abbildung $\mathcal C$.
		\item[(c)] Wie lautet die darstellende Matrix der linearen Abbildung $\mathcal C\circ \mathcal A$?
	\end{itemize}

\vspace{0.3cm}
\noindent\textbf{L�sung:}\newline
\begin{itemize}
	\item[(a)] Die darstellenden Matrizen sind 
		\[ A = \left( \begin{array}{cc} 0 & 0 \\ 3 & -1 \\ 0 & 2 \end{array} \right), \quad B=\left( \begin{array}{cc} -1 & 1 \\ 0 & -1 \\ 0 & 3 \end{array}\right) \,. \]
	\item[(b)] Die Abbildung ist durch \[ \mathcal C:\mathbb R^3 \to \mathbb R^2, \quad \left( \begin{array}{c} x_1 \\x_2 \\x_3\end{array}\right) \mapsto \left( \begin{array}{ccc} -1 & 2 & 0 \\ 0 & 3 & -1 \end{array}\right) \left( \begin{array}{c} x_1 \\x_2 \\x_3\end{array}\right) = \left( \begin{array}{c} -x_1 + 2 x_2 \\3x_2 - x_3 \end{array}\right) \] gegeben. 
	\item[(c)] Es gilt \[ M_{\mathcal C \circ  \mathcal B} = CB= \left( \begin{array}{ccc} -1 & 2 & 0 \\ 0 & 3 & -1 \end{array}\right) \left( \begin{array}{cc} -1 & 1 \\ 0 & -1 \\ 0 & 3 \end{array}\right) =  \left( \begin{array}{cc} 6 & -2 \\ 9 & -5 \end{array}\right)\,.\]
\end{itemize}
\end{exercisebox}

\begin{exercisebox}[Lineare Abbildungen]
	Gegeben sei der reelle Vektorraum 
	\[ \mathcal P_3:= \left\{ p:\mathbb R\to \mathbb R\;|\;p \;\textnormal{Polynome vom Grad } 3\right\}\] und die Abbildung \[ f: \mathcal P_3 \to \mathbb R^2, \quad p\mapsto \left( \begin{array}{c}p(2) \\ p(3)\end{array}\right) \,.\]
	\begin{itemize}
		\item[(a)] Zeigen Sie, dass $f$ eine lineare Abbildung ist und berechnen Sie $f$ f�r die Monome $p_i: x\mapsto x^i,\; i=0,1,2,3\,.$
		\item[(b)] Berechnen Sie $f(p)$ f�r ein allgemeines Polynom $p \in \mathcal P_4$. % mit  \[ p=a_3\cdot p_3 + a_2 \cdot p_2 + a_1\cdot p_1 + a_0 \cdot p_0, \quad a_0,\dots,a_3\in \mathbb R.\] 
		\item[(d)] Zeigen Sie, dass die Funktion \[ g:\mathbb R^4 \to \mathbb R^2, \quad \left( \begin{array}{c}a_3 \\ a_2 \\ a_1 \\ a_0\end{array}\right) \mapsto f(a_3\cdot p_3 + a_2 \cdot p_2 + a_1\cdot p_1 + a_0 \cdot p_0)\] linear ist und berechnen Sie die darstellende Matrix von $g$.
	\end{itemize}

\vspace{0.3cm}
\noindent\textbf{L�sung:}\newline
\begin{itemize}
	\item[(a)] Es gilt f�r $p, q \in \Pi_3$ und $\lambda \in \mathbb R$
		\begin{eqnarray*}
			f(p + q) &=& \left( \begin{array}{c} (p+q) (2) \\ (p + q)(3) \end{array}\right) =  \left( \begin{array}{c} p(2) + q(2) \\ p(3) + q(3) \end{array}\right) =\left( \begin{array}{c} p(2)\\ p(3) \end{array}\right) + \left( \begin{array}{c} q(2) \\ q(3) \end{array}\right) =  f(p) + f(q)\,, \\[1ex]
	f(\lambda \cdot p) &=& \left( \begin{array}{c} \lambda \cdot p(2) \\ \lambda \cdot p(3) \end{array}\right) =  \lambda \cdot \left( \begin{array}{c} p(2) \\ p(3) \end{array}\right) = \lambda \cdot f(p) \,. 
\end{eqnarray*}
Es gilt 
	\[ 
	f(p_0) = \left( \begin{array}{c} 1 \\1 \end{array}\right), \; 
	f(p_1) = \left( \begin{array}{c} 2 \\3 \end{array}\right), \; 
	f(p_2) = \left( \begin{array}{c} 4 \\9 \end{array}\right), \; 
	f(p_3) = \left( \begin{array}{c} 8 \\27 \end{array}\right)\,.
 \]
\item[(b)] Durch die Linearit�t folgt 
	\begin{eqnarray*} 
		f(p) &=& f(a_3\cdot p_3 + a_2\cdot p_2 + a_1\cdot p_1 + a_0\cdot p_0) \\
	&=& a_3  f(p_3) + a_2 f(p_2) + a_1 f(p_1) + a_0f(p_0) \\
	&\stackrel{(b)}{=}&  \left( \begin{array}{c} 8a_3 + 4a_2 + 2a_1 + a_0 \\ 27a_3 + 9a_2 + 3a_1 + a_0 \end{array}\right) 
\end{eqnarray*}
\item[(c)] Es gilt 
	\[ g: \mathbb R^4 \to \mathbb R^2, \quad \left( \begin{array}{c} a_3 \\ a_2 \\ a_1 \\ a_0 \end{array}\right) \mapsto \left( \begin{array}{c} 8a_3 + 4a_2 + 2a_1 + a_0 \\ 27a_3 + 9a_2 + 3a_1 + a_0 \end{array}\right)	
	\]
		Damit ist die darstellende Matrix durch 
		\[ M_g = \left( \begin{array}{cccc} 8 & 4 & 2 &1 \\ 27 & 9 & 3 & 1 \end{array}\right) \]
		gegeben und es gilt 
		\[ g( \boldsymbol a ) = \left( \begin{array}{cccc} 8 & 4 & 2 &1 \\ 27 & 9 & 3 & 1 \end{array}\right) \boldsymbol a \,. \]
\end{itemize}

\end{exercisebox}

 
\newpage

% Lineare Gleichungssysteme
\setcounter{satz}{0} \setcounter{def}{0} \setcounter{exercise}{0}  
	
\ifintegriereaufgaben
\section*{Lineare Gleichungssysteme Aufgaben}\label{MA1-lineare-gleichungssysteme-aufgaben}

	Regeln und Beispiele, die zur Bearbeitung der Aufgaben hilfreich sind, finden Sie im Skript\wennreferenzvorhanden{MA1-lineare-gleichungssysteme}{ (Abschnitt \abschnittsverweis{MA1-lineare-gleichungssysteme})} und, falls Sie nicht weiterkommen, schauen Sie \hyperref[MA1-lineare-gleichungssysteme-loesungen]{hier}.
\else
\section{Lineare Gleichungssysteme Aufgaben}\label{MA1-lineare-gleichungssysteme-aufgaben}

	Regeln und Beispiele, die zur Bearbeitung der Aufgaben hilfreich sind, finden Sie im Skript und, falls Sie nicht weiterkommen, schauen Sie \hyperref[MA1-lineare-gleichungssysteme-loesungen]{hier}.

\fi 

\begin{exercisebox}[Rechenregeln]
	Wie muss $\boldsymbol b$ gew�hlt werden, dass gilt $A\boldsymbol x =\boldsymbol b$ mit \[ \left( \begin{array}{ccccc} -1 & 0 & 0 & 3 & -2 \\  2 & -1 & 1 & 0 & 0 \\ 1 & 1 & -1 & -2 & 0 \end{array}\right), \quad \boldsymbol x = \left( -1, 2,-2,3,5\right)^\intercal\,.\]
\end{exercisebox}

\begin{exercisebox}[L�sbarkeitsanalyse]
	Gegeben seien \[ A = \left( \begin{array}{cc} 1 & 1 \\ 2 & \alpha \end{array}\right), \quad \boldsymbol b = \left( \begin{array}{c} -2 \\ \beta \end{array}\right)\] mit Parametern $\alpha, \beta\in \mathbb R$. 
	\begin{itemize}
		\item[(a)] Bestimmen Sie den Rang der Matrix $A$ in Abh�ngigkeit von $\alpha$.
		\item[(b)] Bestimmen Sie den Rang der erweiterten Matrix $(A|b)$ in Abh�ngigkeit von $\alpha$ und $\beta$.
		\item[(c)] Bestimmen Sie in Abh�ngigkeit der beiden Parameter, wann das lineare Gleichugnssystem $A\boldsymbol x = \boldsymbol b$ l�sbar ist und wann es genau eine L�sung gibt.
	\end{itemize}
\end{exercisebox}

\begin{exercisebox}[Gleichungssysteme]
	Analysieren Sie die L�sbarkeit der folgenden Gleichungssysteme und bestimmen Sie die L�sungsmenge:
	\begin{eqnarray*}
		\begin{array}{ll} 
			(a) \; \begin{cases}
				x_1 -x_2 - x_3 &= 0\\
				x_1 +x_2 + 3x_3 &= 4\\
				     x_2 + 2x_3 &= 2\\
				     2x_1 + x_3 &= 3\\
    \end{cases} 
   &\quad  
    (b)\; \begin{cases}
	    x_1 + x_2 - x_3 + x_4 - x_5 &= 1\\
    -x_1 + x_2 - 3 x_3 + 3x_5 &= 2\\
    x_2 - 2 x_3 + x_4 - x_5 &= -1\\
    2x_1 + 2 x_3 + x_4 - 2 x_5 &= 1
  \end{cases} \\[3ex]
  (c)\; \begin{cases}
	x_1 + x_2 - x_3 &=1 \\
	2x_1 + 3x_2 + px_3 &= 3 \\
	x_1 + px_2 + 3x_3 &= 2
    \end{cases}   
   &\quad  
	(d)\; \begin{cases}
    3x_1 + x_2-4x_3 &= 2\\
    x_1 + x_2&= 0\\
    4x_1 +5x_3 + x_4 &= 1 \\
    6x_2 +x_3 + 2x_4 &= 1
  \end{cases} \\[3ex]
  (e)\;\begin{cases}
	  (2+i)\cdot z_1 + i\cdot z_3 &= -4+3i \\ 
		(1+3i)\cdot z_1 + (-1-i)\cdot z_2  - z_3 &= -6-2i \\ 
		(4+2i)\cdot z_2 + (3+3i)\cdot z_3 &= -2+4i 
	\end{cases} 
  \end{array}
   &\quad  
  \end{eqnarray*}

  \end{exercisebox}

\begin{exercisebox}[Gleichungssysteme]
	Gesucht ist eine Matrix $X\in \mathbb R^{2\times 2}$ so, dass \[ A X + X A^\intercal = C, \quad A=\left( \begin{array}{cc} 1 & -1 \\ 2 & 2 \end{array}\right), \; C = \left( \begin{array}{cc} 4 & 0 \\ 6 & 4 \end{array}\right) \,. \]
	Bestimmen Sie ein lineares Gleichungssystem f�r die unbekannten Koeffizienten $x_{ij},\; i,j=1,2$ der Matrix $X$ und l�sen Sie es.
  \end{exercisebox}

%\begin{exabox}[�berbestimmtes Gleichungssystem]
%	Bestimme die Normalform einer Ebene, wenn die Ebene durch drei Punkte festgelegt ist. 
%  \begin{eqnarray*}
%	  \boldsymbol w = \left( \begin{array}{c} u_2 v_3 - u_3 v_2 \\ u_3 v_1 - u_1 v_3 \\u_1 v_2 - u_2 v_1\end{array}\right) 
%  \end{eqnarray*}
%\end{exabox}

%\begin{exabox}[Analytische Geometrie]
%	Lineares Gleichungssystem l�sen zur expliziten Ebenengleichung bei gegebenen drei Punkten.  - Sp�ter noch mit Determinante?
%\end{exabox}

%\begin{exabox}[Kreuzprodukt] 
%	Geben Sie alle Vektoren $\boldsymbol b\in\mathbb R^3$ an, so dass folgende Bedingungen gelten: 
%	\[ \boldsymbol e_1 \times \boldsymbol b = \boldsymbol e_3, \quad \boldsymbol b \perp \left( \begin{array}{c} 1 \\1 \\ 1\end{array}\right)\,, \] wobei $\boldsymbol e_i$ den $i$-ten Einheitsvektor bezeichnet. Hinweis: Machen Sie hierf�r den Ansatz $\boldsymbol b = \left( b_1, b_2, b_3\right)^\intercal$ und bestimmen Sie ein lineares Gleichungsystem.
%\end{exercisebox}

%\begin{exercisebox}[Direkte Verfahren: bringe auf Dreiecksform]
%	\begin{itemize}
%		\item Gau�sches Eliminationsverfahren 
%		\item Givens-Rotationen
%		\item Householderverfahren
%		\item Cholesky-Zerlegung (positive definite Matrix)
%	\end{itemize}
%\end{exercisebox}
%
%\begin{exercisebox}[Direkte Verfahren, numerischer Aufwand]
%\end{exercisebox}
%
%\begin{exercisebox}[LU-Zerlegung, numerischer Aufwand]
%\end{exercisebox}
%
%\begin{exercisebox}[LR-Zerlegung, numerischer Aufwand]
%\end{exercisebox}
%
%\begin{exercisebox}[Beispiel Matrix aufstellen]
%\end{exercisebox}
%
%\begin{theorembox}[Iterative Methoden - GC-Verfahren, numerischer Aufwand]
%\end{theorembox}
%
%\begin{theorembox}[Iterative Methoden - GMRES, numerischer Aufwand]
%\end{theorembox}
%
%\begin{theorembox}[Norm einer Matrix]
%\end{theorembox}
%
%\begin{exercisebox}[Konditionszahl]
%\end{exercisebox}
%
%\begin{exercisebox}[Vorkonditionierung]
%\end{exercisebox}
%
%\begin{theorembox}[Parallelisierung von Prozessen]
%\end{theorembox}
%
%\begin{exercisebox}[Mathematik lesen, keine eindeutige L�sung]
%\end{exercisebox}
%
%\begin{theorembox}[Rang - MC]
%	\begin{itemize}
%		\item Der Rang der Matrix $A$ ist die Anzahl der nichttrivialen Zeilen in ihrer Zeilennormalform. 
%		\item Der Rang der Matrix $A$ ist die Anzahl der linear unabh�ngigen Spalten. 
%		\item Der Rang der Matrix $A$ ist die Anzahl der linear unabh�ngigen Zeilen. 
%		\item Der Rang einer Matrix ist die Dimension des Bildes. 
%	\end{itemize}
%\end{theorembox}
%
%\begin{theorembox}[Fredholmsche Alternative - Beweis vorf�hren bitte]
%  Entweder $A\boldsymbol x = b$ ist eindeutig l�sbar f�r alle $\boldsymbol b \in \mathbb R^n$ oder $A\boldsymbol x = \boldsymbol 0$ hat eine nichttriviale L�sung. 
%  \[ \Rightarrow A\boldsymbol x = \boldsymbol b \textnormal{ eindeutig f�r alle } \boldsymbol b\in\mathbb R^n, \textnormal{ d.h. } A\boldsymbol x = 0 \textnormal{ eindeutig l�sbar, d.h. } \mathrm{ker}A = \{\boldsymbol 0\}\]
%  \[ \Leftarrow \mathrm{ker}A = \{\boldsymbol 0\}, \textnormal{ so folgt } \mathrm{def}A = 0 \Rightarrow \mathrm{rk}A = n \Rightarrow \textnormal{ eindeutige L�sung}\,. \]
%\end{theorembox}

 
\newpage
\setcounter{satz}{0} \setcounter{def}{0} \setcounter{exercise}{0}  
\section*{Lineare Gleichungssysteme L�sungen}\label{MA1-lineare-gleichungssysteme-loesungen}

\ifintegriereaufgaben
	Regeln und Beispiele, die zum Verst�ndnis der L�sungswege hilfreich sind, finden Sie im Skript\wennreferenzvorhanden{MA1-lineare-gleichungssysteme}{ (Abschnitt \abschnittsverweis{MA1-lineare-gleichungssysteme})}.
\else
	Regeln und Beispiele, die zum Verst�ndnis der L�sungswege hilfreich sind, finden Sie im Skript.
\fi 


\begin{exercisebox}[Rechenregeln]
Wie muss $\boldsymbol b$ gew�hlt werden, dass gilt $A\boldsymbol x =\boldsymbol b$ mit 
\[ \left( \begin{array}{ccccc} -1 & 0 & 0 & 3 & -2 \\  2 & -1 & 1 & 0 & 0 \\ 1 & 1 & -1 & -2 & 0 \end{array}\right), \quad \boldsymbol x = \left( -1, 2,-2,3,5\right)^\intercal\,.\]

  \hspace{0.3cm}
  \newline
  \textbf{L�sung:}
  Zu berechnen ist die Matrix-Vektor Multiplikation
  \begin{eqnarray*}
	  \left( \begin{array}{ccccc} -1 & 0 & 0 & 3 & -2 \\  2 & -1 & 1 & 0 & 0 \\ 1 & 1 & -1 & -2 & 0 \end{array}\right) \left(\begin{array}{c} -1 \\ 2 \\ -2 \\3 \\5 \end{array} \right) 
%&=& \left( \begin{array}{c} 
% (-1)\cdot (-1) + 0\cdot 2 + 0\cdot (-2) + 3 \cdot 3 + (-2) \cdot 5 \\
% 2\cdot (-1) + (-1)\cdot 2 + 1\cdot (-2) + 0 \cdot 3 + 0 \cdot 5 \\
% 1\cdot (-1) + 1\cdot 2 + (-1)\cdot (-2) + (-2) \cdot 3 + 0 \cdot 5 \\
%  \end{array}\right) \\ 
  &=&  \left( \begin{array}{c} 0 \\ -6 \\ -3\end{array}\right) \,.
  \end{eqnarray*}
\end{exercisebox}

\begin{exercisebox}[L�sbarkeitsanalyse]
	Gegeben seien \[ A = \left( \begin{array}{cc} 1 & 1 \\ 2 & \alpha \end{array}\right), \quad \boldsymbol b = \left( \begin{array}{c} -2 \\ \beta \end{array}\right)\] mit Parametern $\alpha, \beta\in \mathbb R$. 
	\begin{itemize}
		\item[(a)] Bestimmen Sie den Rang der Matrix $A$ in Abh�ngigkeit von $\alpha$.
		\item[(b)] Bestimmen Sie den Rang der erweiterten Matrix $(A|b)$ in Abh�ngigkeit von $\alpha$ und $\beta$.
		\item[(c)] Bestimmen Sie in Abh�ngigkeit der beiden Parameter, wann das lineare Gleichugnssystem $A\boldsymbol x = \boldsymbol b$ l�sbar ist und wann es genau eine L�sung gibt.
	\end{itemize}

  \hspace{0.3cm}
  \newline
  \textbf{L�sung:}
  F�r die Bestimmung von $\mathrm{rk}(A)$ und $\mathrm{rk}(A_{\boldsymbol b}$ bringen wir die erweiterte Koeffizientenmatrix $\mathrm{rk}(A_{\boldsymbol b})$ auf Dreiecksform und lesen die gesuchten Gr��en ab: 
		  \begin{eqnarray*}
			  \left( \begin{array}{cc|c} 1 & 1 & -2 \\ 2 & \alpha & \beta \end{array}\right) \Leftrightarrow 
			  \left( \begin{array}{cc|c} 1 & 1 & -2 \\ 0 & \alpha-2 & \beta+4 \end{array}\right) \,. 
		  \end{eqnarray*}
		  \begin{itemize}
			  \item[(a)] 
				  $\mathrm{rk}(A) = \begin{cases} 1 \Leftrightarrow \alpha = 2\,,\\ 2  \Leftrightarrow \alpha \not= 2\,. \end{cases} $
		  \item[(b)] 
			  $\mathrm{rk}(A_{\boldsymbol b}) = \begin{cases} 1 \Leftrightarrow \alpha = 2 \land \beta = -4\,,\\ 2  \Leftrightarrow \alpha \not= 2 \lor \beta \not=-4\,. \end{cases}$ 
		  \item[(c)]  Um zu sehen f�r welche Werte von $\alpha$ und $\beta$ das Gleichungssystem l�sbar ist, bestimmen wir den Fall, f�r den das Gleichungssystem nicht l�sbar ist und negieren dann die Aussage:
			  \begin{eqnarray*}
				  \textnormal{Gleichungssystem nicht l\"osbar } &\Leftrightarrow& \; \alpha = 2 \land \beta \not= -4 \\ 
				  \Rightarrow \; \textnormal{Gleichungssystem l\"osbar } &\Leftrightarrow& \; \neg \left(  \alpha = 2 \land \beta \not= -4\right) \Leftrightarrow    \alpha \not= 2 \lor \beta = -4
			  \end{eqnarray*}
			  Das lineare Gleichungssystem ist eindeutig l�sbar genau dann, wenn $\mathrm{rk}(A) = \mathrm{rk}(A_b) = 2$, also falls $\alpha \not=2 $ und $\beta\not=-4$.
	  \end{itemize}
\end{exercisebox}
\begin{exercisebox}[Gleichungssysteme]
	Analysieren Sie die L�sbarkeit der folgenden Gleichungssysteme und bestimmen Sie die L�sungsmenge:
	\begin{eqnarray*}
		\begin{array}{lll} 
			(a) \; \begin{cases}
				x_1 -x_2 - x_3 &= 0\\
				x_1 +x_2 + 3x_3 &= 4\\
				     x_2 + 2x_3 &= 2\\
				     2x_1 + x_3 &= 3\\
    \end{cases} 
    &\;\;  
    (b)\; \begin{cases}
	    x_1 + x_2 - x_3 + x_4 - x_5 &= 1\\
    -x_1 + x_2 - 3 x_3 + 3x_5 &= 2\\
    x_2 - 2 x_3 + x_4 - x_5 &= -1\\
    2x_1 + 2 x_3 + x_4 - 2 x_5 &= 1
  \end{cases}
  \;\; (c)\; \begin{cases}
	x_1 + x_2 - x_3 &=1 \\
	2x_1 + 3x_2 + px_3 &= 3 \\
	x_1 + px_2 + 3x_3 &= 2
    \end{cases}    \\[3ex]
	(d)\; \begin{cases}
    3x_1 + x_2-4x_3 &= 2\\
    x_1 + x_2&= 0\\
    4x_1 +5x_3 + x_4 &= 1 \\
    6x_2 +x_3 + 2x_4 &= 1
  \end{cases}
    &\;\;  
  (e)\;\begin{cases}
	  (2+i)\cdot z_1 + i\cdot z_3 &= -4+3i \\ 
		(1+3i)\cdot z_1 + (-1-i)\cdot z_2  - z_3 &= -6-2i \\ 
		(4+2i)\cdot z_2 + (3+3i)\cdot z_3 &= -2+4i 
	\end{cases} 
  \end{array}
  \end{eqnarray*}



  \hspace{0.3cm}
  \newline
  \textbf{L�sung:}
	\begin{itemize}
		\item[(a)]  
			Das Gleichungssystem besteht aus vier Gleichungen f�r drei Unbekannte. Hat man mehr Gleichungen als Unbekannte, so spricht man von einem �berbestimmten Gleichungssystem. Das vorliegende �berbestimmte Gleichungssystem besitzt eine eindeutige L�sung und das liegt daran, dass die Gleichungen linear abh�ngig sind. Bringt man die Matrix auf Trapezform, so pssiert folgendes: 
			\begin{eqnarray*}
				\left( \begin{array}{ccc|c} 1 & -1 & -1 & 0 \\ 1 & 1 & 3 & 4 \\ 0 & 1 & 2 & 2 \\ 2 & 0 & 1 & 3 \end{array}\right)  \Leftrightarrow  
				\left( \begin{array}{ccc|c} 1 & -1 & -1 & 0 \\ 0 & 1 & 2 & 2 \\ 0 & 1 & 2 & 2 \\ 0 & 2 & 3 & 3 \end{array}\right) \Leftrightarrow  
					\left( \begin{array}{ccc|c} 1 & -1 & -1 & 0 \\ 0 & 1 & 2 & 2 \\ \textcolor{red}{0} & \textcolor{red}{0} & \textcolor{red}{0} & \textcolor{red}{0} \\ 0 & 0 & -1 & -1 \end{array}\right)  
			\end{eqnarray*}
			Die dritte Zeile ist eine \textcolor{red}{Nullzeile} und da diese Zeile der Gleichung 
			\[ \textcolor{red}{0}\cdot x_1 + \textcolor{red}{0} \cdot x_2 + \textcolor{red}{0}\cdot x_3 = \textcolor{red}{0}\] entspricht, ist sie f�r alle m�glichen Werte von \(x_1, x_2\) und \(x_3\) bedingungslos erf�llt! Es gen�gt die �brigen drei Gleichungen zur Bestimmung der L�sung heranzuziehen und das sieht gut aus: die vierte Gleichung f�hrt auf eine eindeutige L�sung f�r \(x_3\) 
			\[ (-1) \cdot x_3 = -1 \Leftrightarrow x_3 = 1\] 
			Setzt man \(\textcolor{blue}{x_3 = 1 }\) in die zweite Gleichung ein, so erh�lt man eine eindeutige L�sung f�r \(x_2\): 
			\[ \begin{array}{ c rcl l} &  1 \cdot x_2 + 2\cdot \textcolor{blue}{x_3} &=&  2 & \\ \stackrel{ \displaystyle{ \textcolor{blue}{x_3 = 1} } }{\Leftrightarrow} &  x_2 + 2\cdot \textcolor{blue}{1} &=& 2  & \| - 2 \\  \Leftrightarrow &  x_2 &=& 0 & \end{array}\] 
			Im letzten Schritt des R�ckw�rtseinsetzens, setzt man nun die L�sungen f�r \(\textcolor{blue}{x_3}\) und \(\textcolor{blue}{x_2}\) in die erste Gleichugn ein und erh�lt die eindeutige L�sung f�r \(x_1\): 
			\[ \begin{array}{ c rcl l} &  x_1 - \textcolor{blue}{x_2} - \textcolor{blue}{x_3} &=&  0 & \\  \stackrel{ \displaystyle{ \textcolor{blue}{x_3 = 1}, \textcolor{blue}{ x_2 = 0} } }{\Leftrightarrow} &  x_1 - \textcolor{blue}{0} - \textcolor{blue}{1} &=& 0  &  \\  \Leftrightarrow &  x_1 &=& 1 & \end{array}\] 

			Die eindeutige L�sung des Gleichungssystems lautet also $x_1 = 1, x_2 = 0, x_3 = 1$. �berbestimmte Gleichungssysteme sind im allgemeinen nicht l�sbar. In diesem Fall sind die Zeilen linear abh�ngig und das bedeutet, dass die Bedingung, die in einer Zeile drinsteckt durch die anderen schon abgedeckt ist. Eine Zeile kann man aus dem System herausnehmen. Beim L�sen mit dem Gau�-Algorithmus wird diese Information dadurch transparent, dass eine Nullzeile entsteht. 

		\item[(b)] $x_1 = 2-\lambda, x_2 = 1+2\lambda, \; x_3 = \lambda, x_4=-1, x_5=1$ 
		\item[(c)]  
			\begin{eqnarray*}
				\left( \begin{array}{ccc|c} 1 & 1 & - 1 &1 \\ 2 & 3 & p & 3 \\ 1 & p & 3 & 2 \end{array}\right) {\Leftrightarrow }
				\left( \begin{array}{ccc|c} 1 & 1 & - 1 &1 \\ 0 &1 & 2+p & 1 \\ 0 & p-1 & 4 & 1 \end{array}\right) \stackrel{ p\not=1}{\Leftrightarrow }
				\left( \begin{array}{ccc|c} 1 & 1 & - 1 &1 \\ 0 & 1 & 2+p & 1 \\ 0 & 0 & 4 + (1-p)(2+p) & 2-p \end{array}\right)\,, 
			\end{eqnarray*}
			wobei die faktorisierte Form des Eintrags an Position mit Index \(33\) lautet: 
			\[ 4+(1-p)(2+p) = -p^2-p+6 = -(x-2)(x+3)\,.\]
			Das hei�t, dass die F�lle \( p=1\), \(p=2\), \(p=-3\) und \(p\in \mathbb R\setminus\{1,2,-3\}\) separat zu untersuchen sind: 
			\begin{enumerate}
				\item F�r $p=1$ lautet die L�sung $x_1=1, x_2 = x_3 = \displaystyle{\frac14}\,.$  
				\item F�r $p=-3$ ist das System nicht l�sbar.  
				\item F�r $p=2$ hat das System unendlich viele L�sungen: $x_1 5\lambda, \, x_2 = 1-4\lambda, \; x_3 = \lambda, \; \lambda \in \mathbb R\,.$  
				\item F�r $p\in \mathbb R\setminus\{1,-3,2\}$ lautet die eindeutige L�sung $x_1 = 1, x_2 = x_3 =\displaystyle{ \frac{1}{3+p} }\,.$
			\end{enumerate}
			
  \item[(d)]  
			\begin{eqnarray*}
				\left( \begin{array}{cccc|c} 3 & 1 & -4 & 0 & 2 \\ 1 & 1 & 0 & 0 & 0 \\ 4 & 0 & 5 & 1 & 1 \\ 0 & 6 & 1 & 2 & 1 \end{array}\right) &\Leftrightarrow & 
				\left( \begin{array}{cccc|c}  1 & 1 & 0 & 0 & 0 \\ 3 & 1 & -4 & 0 & 2 \\4 & 0 & 5 & 1 & 1 \\ 0 & 6 & 1 & 2 & 1 \end{array}\right) \\
				&\Leftrightarrow & 
				\left( \begin{array}{cccc|c}  1 & 1 & 0 & 0 & 0 \\ 0 & 1 & 2 & 0 & -1 \\0 & 0 & 2 & 3 & 4 \\ 0 & 0 & 0 & 37 & 58 \end{array}\right)  \Rightarrow \left( \begin{array} {c} x_1 \\x_2 \\x_3 \\ x_4 \end{array}\right) =  \frac{1}{37}\left( \begin{array} {c} 41 \\-41 \\ 2 \\ 58 \end{array}\right) 
			\end{eqnarray*}

  \item[(e)] 
			\begin{eqnarray*}
				\left( \begin{array}{ccc|c} 2+i & 0 & j & -4+3i \\ 1+3i & -1-i & -1 & -6 -2i \\ 0 & 4+2i & 3+3i & -2 + 4i \end{array}\right) &\Leftrightarrow & 
				\left( \begin{array}{ccc|c} 2+i & 0 & j & -4+3i \\ 0 & -1-i & -i & -6 1-i \\ 0 & 4+2i & 3+3i & -2 + 4i \end{array}\right) \\ &\Leftrightarrow  &
				\left( \begin{array}{ccc|c} 2+i & 0 & j & -4+3i \\ 0 & -1-i & -i & -6 1-i \\ 0 & 0 & 2 & 0 \end{array}\right)
			\end{eqnarray*}
			L�sen wir die Gleichungen von unten nach oben auf, dann erh�lten wir 
			\[ z_1 = -1+2i, \; z_2 = i, \; z_3 = 0\,.\]
	  \end{itemize}
\end{exercisebox}
\begin{exercisebox}[Gleichungssysteme]
	Gesucht ist eine Matrix $X\in \mathbb R^{2\times 2}$ so, dass \[ A X + X A^\intercal = C, \quad A=\left( \begin{array}{cc} 1 & -1 \\ 2 & 2 \end{array}\right), \; C = \left( \begin{array}{cc} 4 & 0 \\ 6 & 4 \end{array}\right) \,. \]
	Bestimmen Sie ein lineares Gleichungssystem f�r die unbekannten Koeffizienten $x_{ij},\; i,j=1,2$ der Matrix $X$ und l�sen Sie es.

  \hspace{0.3cm}
  \newline
  \textbf{L�sung:}
  Es muss gelten 
  \[ 
  \left( \begin{array}{cc} -x_{12} - x_{21} + 2 x_{11} & -x_{22}+3x_{12} + 2x_{11} \\ -x_{22} + 3 x_{21} + 2 x_{11}  & 4 x_{22} + 2 x_{12} + 2 x_{21}  \end{array}\right) = \left( \begin{array}{cc} 4 & 0 \\6 & 4 \end{array}\right) \,. \]
  Also 
  \[ 
  \left( \begin{array}{cccc| c } 2 & -1 & -1 & 0 & 4 \\ 2 & 0 & 3 & -1 & 0 \\ 2 & 3 & 0 & -1 &  6 \\ 0 & 2 & 2 & 4 & 4 \end{array}\right) \]
  Dieses System l�sen wir mittels Gau�-Verfahren und erhalten 
  \[ \Leftrightarrow 
  \left( \begin{array}{cccc| c } 2 & -1 & -1 & 0 & 4 \\ 0 & 1 & 4 & -1 & 4 \\ 0 & 0 & -5 & 1 &  6 \\ 0 & 0 & 0 & 1 & 1 \end{array}\right) \Rightarrow x_{22}= 1, x_{21} = -1, x_{12} = 1, x_{11} = 2\,. \]
  \end{exercisebox}
 
\newpage

% Inverse
\setcounter{satz}{0} \setcounter{def}{0} \setcounter{exercise}{0}  

\ifintegriereaufgaben
\section*{Determinante, inverse Matrix Aufgaben}\label{MA1-determinante-und-inverse-aufgaben}

	Regeln und Beispiele, die zur Bearbeitung der Aufgaben hilfreich sind, finden Sie im Skript\wennreferenzvorhanden{MA1-determinante-und-inverse}{ (Abschnitt \abschnittsverweis{MA1-determinante-und-inverse})} und, falls Sie nicht weiterkommen, schauen Sie \hyperref[MA1-determinante-und-inverse-loesungen]{hier}.
\else
\section{Determinante, inverse Matrix Aufgaben}\label{MA1-determinante-und-inverse-aufgaben}

	Regeln und Beispiele, die zur Bearbeitung der Aufgaben hilfreich sind, finden Sie im Skript und, falls Sie nicht weiterkommen, schauen Sie \hyperref[MA1-determinante-und-inverse-loesungen]{hier}.
\fi 

\begin{exercisebox}[Determinante]
	Berechnen Sie die Determinante der Matrix \[ A = \left( \begin{array}{ccc} 2 & -6 & 2 \\ 0 & 0 & 2 \\ -1 & 2 & 1 \end{array}\right)\,.  \]

\end{exercisebox}

\begin{exercisebox}[Determinante]
	Betrachten Sie die Schar der Matrizen \[ A_a = \left( \begin{array}{ccc} 1 & 1 & a \\ 1 & a & 1 \\ 1 & a & a \end{array}\right)\,, \quad a \in \mathbb R\,. \] 
	\begin{itemize}
		\item[(a)] Berechnen Sie die Determinante von $A_a$.
		\item[(b)] F�r welche $a\in \mathbb R$ ist $A_a$ regul�r?
		\item[(c)] Berechnen Sie mit Hilfe des Gau�-Algorithmus die inversen Matrizen $A_a^{-1}$ f�r alle regul�ren Matrizen $A_a$.
		\item[(d)] Berechnen Sie mit Hilfe der Cramerschen Regel die erste Spalte der inversen Matrizen $A_a^{-1}$ f�r alle regul�ren Matrizen $A_a$. Hinweis: L�sen Sie $ A_a \boldsymbol v = \boldsymbol e_1, \quad \boldsymbol e_1 = (1, 0, 0)^\intercal\,.$

\end{itemize}

\end{exercisebox}
\begin{exercisebox}[Basis]
	�berpr�fen Sie mit Hilfe der Determinante, ob die Vektoren 
	\[ 
		\boldsymbol v_1 = ( 2, -1, 0)^\intercal, \quad  
		\boldsymbol v_2 = ( 1, 1, 2 )^\intercal, \quad 
	\boldsymbol v_3 = ( -5, 1,-2 )^\intercal \]
	eine Basis des $\mathbb R^3$ sind.
\end{exercisebox}

\begin{exercisebox}[Gleichungssysteme]
	\begin{itemize}
		\item[(a)] Zeigen Sie, dass die Matrix 
			\[ A = \left( \begin{array}{ccc} 1 & 2 & 0 \\ 1 & 7 & 4 \\ 3 & 13 & 4 \end{array}\right) \]
			regul�r ist. 

		\item[(b)] L�sen Sie das lineare Gleichungssystem f�r $x, y, z \in \mathbb R$
	\[ \begin{array}{ccccccc} 
			x_1 & + & 2x_2 & && = & x \\ 
			x_1 & + & 7x_2 &+ & 4x_3 & = & y \\ 
			3x_1 & + & 13x_2 &+ & 4x_3 & = & z 
	\end{array} \]
		\item[(c)] L�sen Sie das lineare Gleichungssystem f�r $x=3, y=8, z=0$
	\end{itemize}
\end{exercisebox}

\begin{exercisebox}[Anwendungen: Eigenwerte]
	Berechnen Sie alle $\lambda \in \mathbb C$, f�r die die Matrix $A- \lambda I_2$ nicht regul�r ist, wobei \[ A = \left( \begin{array}{cc} -11 & 12 \\ -9 & 10 \end{array}\right)\]. 
\end{exercisebox}

%\begin{itemize}
%	\item Erkl�rung: antisymmetrische Multilinearform
%	\item Spatprodukt, Vektorprodukt, Skalarprodukt
%	\item Kleinste Fehlerquadrate - lineaeres Gleichungssystem
%\end{itemize}

%\subsection{Details und Extras}
%\begin{mybox}[Inverse Matrix]
%  Es seien $A,B \in \mathbb R^{n\times n}$ invertierbar. Es gilt 
%  \begin{itemize}
%    \item $\left( A^{-1}\right)^{-1} = A$
%    \item $\left( A B\right)^{-1} = B^{-1} A^{-1}$
%  \end{itemize}
%\end{mybox}
 
\newpage
\setcounter{satz}{0} \setcounter{def}{0} \setcounter{exercise}{0}  
\section*{Determinante, inverse Matrix L�sungen}\label{MA1-determinante-und-inverse-loesungen}

\ifintegriereaufgaben
	Regeln und Beispiele, die zum Verst�ndnis der L�sungswege hilfreich sind, finden Sie im Skript\wennreferenzvorhanden{MA1-determinante-und-inverse}{ (Abschnitt \abschnittsverweis{MA1-determinante-und-inverse})}.
\else
	Regeln und Beispiele, die zum Verst�ndnis der L�sungswege hilfreich sind, finden Sie im Skript.
\fi 

\begin{exercisebox}[Determinante]
	Berechnen Sie die Determinante der Matrix \[ A = \left( \begin{array}{ccc} 2 & -6 & 2 \\ 0 & 0 & 2 \\ -1 & 2 & 1 \end{array}\right)\,.  \]

\vspace{0.3cm}
\noindent\textbf{L�sung:}\newline
Man kann die Determinante auf unterschiedliche Weisen berechnen: 
\begin{itemize}
	\item[(a)] Entwicklungssatz nach der ersten Spalte (Entwicklungssatz nach Laplace)
		\[ \mathrm{det}(A) = 2 \cdot ( 0\cdot 1 - 2 \cdot 2 ) - 0 \cdot ( -6 \cdot 1 - 2 \cdot 2 ) + (-1) \cdot (-6 \cdot 2 - 0 \cdot 2 ) = 4 \,. \]
	\item[(b)] Regel von Sarrus
		\[ \mathrm{det}(A) =  { \color{green}+2 \cdot 0 \cdot 1 + (-6)\cdot 2 \cdot (-1) + 2 \cdot 0\cdot 2}  {\color{blue}- (-1)\cdot 0 \cdot 2 - 2 \cdot 2 \cdot 2 - 1 \cdot 0 \cdot (-6)} = 12 - 8 = 4  \,. \]
\end{itemize}

\textbf{Bemerkung:} Hat eine Matrix Dreiecksgestalt, dann ist die Determinante gerade das Produkt der Diagonalelemente. Also k�nnte man die Determinante der gegebenen Matrix auch berechnen, indem man die Matrix auf Dreiecksgestalt bringt. Aber Vorsicht: Vertauschen von Zeilen und skalare Vielfache ver�ndern den Wert der Determinante und m�ssen ber�cksichtigt werden!    
\begin{eqnarray*} 
	\mathrm{det}(A) &=&  \left| \begin{array}{ccc} 2 & -6 & 2 \\  0 & 0 & 2 \\ -1 & 2 & 1 \end{array}  \right| = {\color{red}-} \left| \begin{array}{ccc} 2 & -6 & 2 \\ {\color{red} -1} & {\color{red}2} & {\color{red} 1} \\ {\color{red} 0 } &  {\color{red} 0} & {\color{red} 2} \end{array}  \right|  = - {\color{green}2}\cdot {\color{blue}2} \left| \begin{array}{ccc} {\color{green}1} &  {\color{green}-3} & {\color{green}1} \\ -1 & 2 & 1 \\ 0 & 0 & {\color{blue}1} \end{array}  \right| =  - 4 \left| \begin{array}{ccc} 1 & -3 & 1 \\ 0 & -1 & 2 \\ 0 & 0 & 1 \end{array}  \right| \\  &=& -4 \cdot ( 1 \cdot (-1)\cdot 1 ) = 4\,.
\end{eqnarray*}
\end{exercisebox}

\begin{exercisebox}[Determinante]
	Betrachten Sie die Schar der Matrizen \[ A_a = \left( \begin{array}{ccc} 1 & 1 & a \\ 1 & a & 1 \\ 1 & a & a \end{array}\right)\,, \quad a \in \mathbb R\,. \] 
	\begin{itemize}
		\item[(a)] Berechnen Sie die Determinante von $A_a$.
		\item[(b)] F�r welche $a\in \mathbb R$ ist $A_a$ regul�r?
		\item[(c)] Berechnen Sie mit Hilfe des Gau�-Algorithmus die inversen Matrizen $A_a^{-1}$ f�r alle regul�ren Matrizen $A_a$.
		\item[(d)] Berechnen Sie mit Hilfe der Cramerschen Regel die erste Spalte der inversen Matrizen $A_a^{-1}$ f�r alle regul�ren Matrizen $A_a$. Hinweis: L�sen Sie $ A_a \boldsymbol v = \boldsymbol e_1, \quad \boldsymbol e_1 = (1, 0, 0)^\intercal\,.$

\end{itemize}

\vspace{0.3cm}
\noindent\textbf{L�sung:}\newline
\begin{itemize}
	\item[(a)] $\mathrm{det}(A_a) = 1 \cdot ( a\cdot a -a\cdot 1) - 1 \cdot (1\cdot a-a\cdot a) + 1 \cdot (1 \cdot 1 - a\cdot a) =  a^2 - 2 a +1 = (a-1)^2\,.$
	\item[(b)] $\mathrm{det}(A_a) \not= 0 \Leftrightarrow a\not=1 \; \Rightarrow $ $A_a$ ist regul�r f�r alle $a\in \mathbb R \setminus\{ 1\}\,.$
	\item[(c)] F�r $a\not=1$ l��t sich die Inverse zu $A_a$ bestimmen durch 
		\[ \left( A_a | I_3 \right) \stackrel{\textnormal{�ZU}}{ \Rightarrow } \; \left( I_3 | A_a^{-1}\right) \,.\]
		Das hei�t man betrachtet das um die Einheitsmatrix erweiterte System  $ ( A_a | I_3)$ und transformiert die linke Seite $A_a$ durch �quivalente Zeilenumformungen (�ZU) zu $I_3$. Aus der rechten Seite $I_3$ wird dann simultan die inverse Matrix $A_a^{-1}$:
		\begin{eqnarray*}
		&&	\left( \begin{array}{ccc|ccc} 1 & 1 & a & 1 & 0 & 0 \\ 1 & a& 1 & 0 & 1 & 0 \\  1 & a & a & 0 & 0 &1 \end{array}\right) \Leftrightarrow
		 \left( \begin{array}{ccc|ccc} 1 & 1 & a & 1 & 0 & 0 \\ 0 & 1-a& a-1 & 1 & -1 & 0 \\  0 & 1-a & 0 & 1 & 0 & -1 \end{array}\right) \\
		 &\Leftrightarrow& 
		 \left( \begin{array}{ccc|ccc} 1 & 1 & a & 1 & 0 & 0 \\ 0 & 1-a& a-1 & 1 & -1 & 0 \\  0 & 0 & 1-a & 0 & 1 & -1 \end{array}\right)  
		 \stackrel{a\not\in\{1,0\}}{\Leftrightarrow}
		 \left( \begin{array}{ccc|ccc} a-1 & a-1 & 0 & a-1 & a & -a \\ 0 & 1-a & 0 & 1 & 0 & -1 \\  0 & 0 & 1-a & 0 & 1 & -1 \end{array}\right) \\
		 &\Leftrightarrow& 
		 \left( \begin{array}{ccc|ccc} a-1 & 0 & 0 & a & a & -a-1 \\ 0 & 1-a & 0 & 1 & 0 & -1 \\  0 & 0 & 1-a & 0 & 1 & -1 \end{array}\right) 
		 \Leftrightarrow
		 \left( \begin{array}{ccc|ccc} 1 & 0 & 0 & \displaystyle{\frac{a}{a-1}} & \displaystyle{\frac{a}{a-1}} & \displaystyle{\frac{-a-1}{a-1}} \\ 0 & 1 & 0 & \displaystyle{\frac{1}{a-1}} & 0 & \displaystyle{\frac{-1}{a-1}} \\  0 & 0 & 1 & 0 & \displaystyle{\frac{1}{1-a}} & \displaystyle{\frac{-1}{1-a}} \end{array}\right) 
	 \end{eqnarray*}
	 Die zu $A_a$ inverse Matrix f�r $a\in \mathbb R\setminus \{ 1,0\}$ lautet 
	 \[ A^{-1}_a = \frac{1}{1-a} \left( \begin{array}{ccc} -a & -a & 1+a \\ 1 & 0 & -1 \\ 0 & 1  & - 1 \end{array}\right)\,. \]
	 Dass die Matrix f�r \(a=1\) nicht invertierbar ist, wissen wir. Allerdings sollten wir eine inverse Matrix bestimmen k�nnen f�r \(a=0\). Die soeben berechnete Matrix \(A^{-1}_a\) ist auch f�r \(a=0\) wohldefiniert und das legt die Vermutung nah, dass \(A^{-1}_0\) die inverse zu \(A_0\) ist.�berpr�fen wir das: 
	 \[ A^{-1}_0 A = \left( \begin{array}{ccc} 0 & 0 & 1 \\ 1 & 0 & -1 \\ 0 & 1  & - 1 \end{array}\right)  \left( \begin{array}{ccc} 1 & 1 & 0 \\ 1 & 0 & 1 \\ 1 & a & 0 \end{array}\right) = \left( \begin{array}{ccc} 1 & 0 & 0 \\ 0 & 1 & 0 \\ 0 & 0 & 1 \end{array}\right)\]
	 Also ist \(A^{-1}_a\) die inverse zu \(A\) f�r alle \(a\in \mathbb R \setminus \{ 1\}\).

 \item[(c)] Die Cramersche Regel kann angewand werden, um eine L�sungskomponente $x_i$ eines Gleichungssystems $A\boldsymbol x = \boldsymbol b$ zu berechnen, es gilt 
	 \[ x_i = \frac{1}{\mathrm{det}(A)} \cdot \mathrm{det}\left( \boldsymbol a_1 \, \vdots  \,\boldsymbol a_2  \,\vdots  \,\boldsymbol a_3  \,\vdots  \,\cdots  \,\vdots   \,\boldsymbol a_{i-1}  \,\vdots   \,\boldsymbol b  \,\vdots  \,\boldsymbol a_{i+1}  \,\vdots \, \cdots \, \vdots  \,\boldsymbol a_n \right)\,. \]
	 Das hei�t im konkreten Fall
	 \begin{eqnarray*} 
		 v_1 &=& \frac{1}{(a-1)^2} \cdot \mathrm{det}\left( \begin{array}{ccc}  1 & 1 & a \\ 0 & a & 1 \\ 0 & a & a \end{array}\right) =  \frac{1}{(a-1)^2} \left( 1\cdot( a\cdot a - a\cdot 1 )\right) = \frac{a}{a-1}\,, \\
		 v_2 &=& \frac{1}{(a-1)^2} \cdot \mathrm{det}\left( \begin{array}{ccc}  1 & 1 & a \\ 1 & 0 & 1 \\ 1 & 0 & a \end{array}\right) =  \frac{1}{(a-1)^2}\left( - (1\cdot a - 1 \cdot 1) \right) =  \frac{1}{1-a}\,,   \\
		 v_3 &=& \frac{1}{(a-1)^2} \cdot \mathrm{det}\left( \begin{array}{ccc}  1 & 1 & 1 \\ 1 & a & 0 \\ 1 & a & 0 \end{array}\right) \frac{1}{(a-1)^2} \left( 1 \cdot (1 \cdot a - 1 \cdot a ) \right) = 0\,,
	 \end{eqnarray*}
	 also die erste Spalte der inversen Matrix $A_a^{-1}$ f�r $a\not=1$. 
\end{itemize}
\end{exercisebox}
\begin{exercisebox}[Basis]
	�berpr�fen Sie mit Hilfe der Determinante, ob die Vektoren 
	\[ 
		\boldsymbol v_1 = ( 2, -1, 0)^\intercal, \quad  
		\boldsymbol v_2 = ( 1, 1, 2 )^\intercal, \quad 
	\boldsymbol v_3 = ( -5, 1,-2 )^\intercal \]
	eine Basis des $\mathbb R^3$ sind.
	
\vspace{0.3cm}
\noindent\textbf{L�sung:}\newline
Die Vektoren sind genau dann eine Basis, wenn jeder Vektor $\boldsymbol w\in \mathbb R^3$ eindeutig als eine lineare Kombination der Vektoren $\boldsymbol v_i, \; i=1,2,3$ \[
\boldsymbol w = \lambda_1 \boldsymbol v_1 + \lambda_2 \boldsymbol v_2 +   \lambda_3 \boldsymbol v_3\]
geschrieben werden kann, d.h. wenn das lineare Gleichungssystem 
\[ V \boldsymbol \lambda = \boldsymbol w \quad \textnormal{mit} \; 
	V= \left( \boldsymbol v_1 \,\vdots \, \boldsymbol v_2  \, \vdots \, \boldsymbol v_3\right), \quad 
\boldsymbol \lambda = (\lambda_1, \lambda_2, \lambda_3 )^\intercal \,, \]
f�r jede rechte Seite $\boldsymbol w \in \mathbb R^3$ eindeutig l�sbar ist. Das ist genau dann der Fall, wenn die Matrix $V \in \mathbb R^{3\times 3}$ mit Spaltenvektoren $\boldsymbol v_i$, $i=1,2,3$ regul�r ist. Also berechnen wir die Determinante
\[ \mathrm{det}(V) = 2\cdot ( 1\cdot(-2) - 2 \cdot 1 ) - (-1) \cdot (1\cdot(-2) - 2 \cdot (-5) ) + 0 \cdot ( 1\cdot 1 - 1 \cdot (-5) ) = -8 + 8 = 0\,, \] 
und stellen fest, dass Matrix nicht regul�r ist und die Vektoren damit keine Basis bilden.
\end{exercisebox}

\begin{exercisebox}[Gleichungssysteme]
	\begin{itemize}
		\item[(a)] Zeigen Sie, dass die Matrix 
			\[ A = \left( \begin{array}{ccc} 1 & 2 & 0 \\ 1 & 7 & 4 \\ 3 & 13 & 4 \end{array}\right) \]
			regul�r. 

		\item[(b)] L�sen Sie das lineare Gleichungssystem f�r $x, y, z \in \mathbb R$
	\[ \begin{array}{ccccccc} 
			x_1 & + & 2x_2 & && = & x \\ 
			x_1 & + & 7x_2 &+ & 4x_3 & = & y \\ 
			3x_1 & + & 13x_2 &+ & 4x_3 & = & z 
	\end{array} \]
		\item[(c)] L�sen Sie das lineare Gleichungssystem f�r $x=3, y=8, z=0$
	\end{itemize}

	
\vspace{0.3cm}
\noindent\textbf{L�sung:}\newline
\begin{itemize}
	\item[(a)] $\mathrm{det}(A) = -4\cdot ( 1\dot 13 - 3 \cdot 2 ) + 4 \cdot ( 1 \dot 7 - 1 \cdot 2 ) = -4 \cdot 7 + 4 \cdot 5 = -8 \not=0 \Rightarrow $ $A$ ist regul�r.
	\item[(b)] Ist die Matrix $A$ einer linearen Abbildung $\mathcal A$ regul�r, dann l��t sich die Abbildung invertieren und die inverse  Abbildung $\mathcal A^{-1}$ ist wieder eine lineare Abbildung. Die Matrix, die $\mathcal A^{-1}$ darstellt ist die zu $A$ inverse Matrix $A^{-1}$. Wenn man ein Gleichungssystem f�r beliebige rechte Seiten l�sen soll, dann bietet es sich an $A^{-1}$ einmal zu berechnen. F�r eine gegebene rechte Seite $\boldsymbol b$ berechnet sich dann die L�sung des entsprechenden Gleichungssystems durch die Matrix-Vektor-Multiplikation  $A^{-1}\boldsymbol b$.
		\begin{eqnarray*} 
		&&	\left( A \,| \, I_3\right)  = 
		\left( \begin{array}{ccc | ccc } 1 & 2 & 0 & 1 0 & 0 \\ 1 & 7 & 4 & 0 &1 & 0 \\ 3 & 13 & 4 & 0 & 0 & 1 \end{array}\right) 
		\Leftrightarrow 
		\left( \begin{array}{ccc | ccc } 1 & 2 & 0 & 1 0 & 0 \\ 0 & 5 & 4 & -1 & 1 & 0 \\ 0 & 7 & 4 & -3 & 0 & 1 \end{array}\right)   \\
		&&\Leftrightarrow 
		\left( \begin{array}{ccc | ccc } 1 & 2 & 0 & 1 0 & 0 \\ 0 & 5 & 4 & -1 & 1 & 0 \\ 0 & 0 & -8 & -8 & -7 & 5 \end{array}\right)  
		 \Leftrightarrow 
		\left( \begin{array}{ccc | ccc } 1 & 2 & 0 & 1 0 & 0 \\ 0 & 10 & 0 & -10 & -5 & 5 \\ 0 & 0 & -8 & -8 & -7 & 5 \end{array}\right)  \\
		 && \Leftrightarrow  
		\left( \begin{array}{ccc | ccc } -5 & 0 & 0 & -15 & -5 & 5 \\ 0 & 10 & 0 & -10 & -5 & 5 \\ 0 & 0 & -8 & -8 & -7 & 5 \end{array}\right) \quad 
		\Rightarrow A^{-1} = \left( \begin{array}{ccc } 3 & 1 & -1 \\ -1 & -\frac12 & \frac12 \\ 1  & \frac78 & -\frac58 \end{array}\right) 
\end{eqnarray*}
\item[(c)] Wir setzen $\boldsymbol b = (x,y,z)^\intercal$ und berechnen $ A^{-1} \boldsymbol b = ( 11, -7 , 10)^\intercal$.
	\end{itemize}
\end{exercisebox}

\begin{exercisebox}[Anwendungen: Eigenwerte]
	Berechnen Sie alle $\lambda \in \mathbb C$, f�r die die Matrix $A- \lambda I_2$ nicht regul�r ist, wobei \[ A = \left( \begin{array}{cc} -11 & 12 \\ -9 & 10 \end{array}\right)\]. 

\vspace{0.3cm}
\noindent\textbf{L�sung:}\newline
Zu bestimmen sind die sogenannten Eigenwerte der Matrix $A$. Sie sind L�sungen der Gleichung 
\begin{eqnarray*}
	\mathrm{det}( A-\lambda I_2) = 0\,.
\end{eqnarray*}
Also berechnen wir die Determinante in Abh�ngigkeit von $\lambda$ und l�sen die resultierende quadratische Gleichung:
\begin{eqnarray*}
	\mathrm{det}\left( \begin{array}{cc} -11 - \lambda & 12 \\ -9 & 10 - \lambda \end{array}\right) &=& (-11-\lambda) \cdot (10-\lambda) - 12\cdot (- 9) =  \lambda^2 +\lambda - 2 \\
	&&\Rightarrow \mathrm{det}( A-\lambda I_2) = 0\Leftrightarrow \lambda_{1,2}=  \frac{-1 \pm \sqrt{1-4\cdot 1 \cdot (-2) }}{2\cdot 1} = \begin{cases} 1 \\ -2 \end{cases}
\end{eqnarray*}
\end{exercisebox}
 

}

\end{document}
